\vspace{1cm}
\begin{flushright}
    \textbf{\LARGE Tóm tắt}
\end{flushright}
\vspace{1cm}
Trong kỷ nguyên số hóa, Internet vạn vật (IoT) đã chuyển mình từ một khái niệm công nghệ trở thành hạ tầng thiết yếu của đời sống hiện đại. Từ các thiết bị dân dụng như Smart TV đến các mạng lưới cảm biến công nghiệp phức tạp, IoT đang tạo ra một lượng dữ liệu khổng lồ. Xu hướng tất yếu hiện nay là tích hợp Trí tuệ nhân tạo (AI) vào các thiết bị này để tạo nên AIoT (Artificial Intelligence of Things), giúp thiết bị không chỉ "kết nối" mà còn có khả năng "tư duy".

Tuy nhiên, việc triển khai AI trên các thiết bị IoT truyền thống đang vấp phải "nút thắt cổ chai" về kiến trúc. Mô hình phổ biến hiện nay dựa vào việc đẩy dữ liệu lên máy chủ đám mây (Cloud Server) để xử lý tập trung. Mô hình này bộc lộ nhiều điểm yếu chí mạng: rủi ro gián đoạn dịch vụ khi mất kết nối mạng, độ trễ cao không phù hợp cho các ứng dụng thời gian thực, và gánh nặng băng thông lớn.

Về mặt phần cứng, các vi xử lý (CPU) thương mại từ các hãng lớn, dù mạnh mẽ, thường gặp vấn đề về tiêu thụ năng lượng khi phải gồng gánh các tác vụ AI liên tục. Hơn nữa, tính chất đóng kín của các dòng chip này khiến các nhà phát triển gặp khó khăn trong việc tùy biến sâu để tối ưu hóa hiệu năng cho các thuật toán AI cụ thể, đồng thời rào cản chi phí bản quyền cũng là một trở ngại lớn cho đổi mới sáng tạo.

Sự xuất hiện của RISC-V đã mang đến một luồng gió mới cho ngành công nghiệp bán dẫn. Là một kiến trúc tập lệnh mở (Open ISA), RISC-V phá bỏ thế độc quyền của các kiến trúc đóng, trao quyền cho các tổ chức nhỏ, các nhóm nghiên cứu và cá nhân khả năng tự thiết kế các lõi CPU chuẩn hóa. Ưu điểm vượt trội của RISC-V nằm ở tính mô-đun hóa và khả năng mở rộng. Các kỹ sư có thể tự do thêm các tập lệnh tùy chỉnh hoặc sửa đổi vi kiến trúc để phục vụ cho các tác vụ chuyên biệt mà không vướng phải rào cản pháp lý hay chi phí bản quyền đắt đỏ. Đây chính là chìa khóa để giải quyết bài toán hiệu năng/năng lượng cho các thiết bị IoT thế hệ mới.

Dựa trên những phân tích trên, dự án của nhóm tập trung phát triển một nền tảng Edge-AI (AI tại biên) chuyên dụng, lấy kiến trúc RISC-V làm nòng cốt. Mục tiêu của nhóm không chỉ dừng lại ở việc sử dụng một lõi CPU đa năng. Thay vào đó, nhóm tận dụng khả năng tùy biến của RISC-V để thiết kế một hệ thống tích hợp sâu, bao gồm một lõi xử lý trung tâm được tối ưu hóa kết hợp với một bộ tăng tốc AI tùy chỉnh
